% Source: Weinberg GR, physical PDF pp. 494--499
%         (printed pp. 470--475; boundary pages are partial).
% Coverage: Section 15.1, Einstein's Equations; equations
%           (15.1.1)--(15.1.26), including all unnumbered displays.
% Figures/tables/footnotes: reference notes 1 and 2; no figures or tables.
% Status: source-reviewed and compile-clean.
% Uncertainties: curvature/sign modernization checks are documented below;
%                no unresolved source readings.
% MODERNIZATION: R_{\mathrm{source}}(t) -> a(t); H=\dot a/a. Compact
% partial derivatives replace long coordinate quotients. Curvature-dependent
% formulas use the binding convention
% R^\rho{}_{\sigma\mu\nu}
% =\partial_\mu\Gamma^\rho{}_{\nu\sigma}
% -\partial_\nu\Gamma^\rho{}_{\mu\sigma}+\cdots, opposite to Weinberg's
% printed convention. The Friedmann equations themselves are unchanged.

\subsection{Einstein's Equations}\label{sec:15.1}

Let us begin our discussion of dynamical cosmology by considering the
constraints imposed by Einstein's field equations on the metric for a general
isotropic and homogeneous universe. According to the results of
Section~\ref{sec:14.2}, this metric can be chosen to have the
Robertson--Walker form:
\begin{equation}
  g_{tt}=-1,\qquad
  g_{it}=0,\qquad
  g_{ij}=a^2(t)\widetilde g_{ij}(\mathbf{x}).
  \tag{15.1.1}\label{eq:15.1.1}
\end{equation}
Here \(t\) is a cosmic time coordinate; \(i\) and \(j\) run over three
comoving spatial coordinates \(r,\theta,\varphi\); and
\(\widetilde g_{ij}\) is the metric for a three-dimensional maximally
symmetric space:
\begin{equation}
  \widetilde g_{rr}=(1-kr^2)^{-1},\qquad
  \widetilde g_{\theta\theta}=r^2,\qquad
  \widetilde g_{\varphi\varphi}=r^2\sin^2\theta,\qquad
  \widetilde g_{ij}=0\quad(i\ne j),
  \tag{15.1.2}\label{eq:15.1.2}
\end{equation}
with \(k\) equal to \(+1\), \(-1\), or \(0\).

The only nonvanishing elements of the affine connection for this metric are
\begin{align}
  \Gamma^t{}_{ij}
  &=a\dot a\,\widetilde g_{ij},
  \tag{15.1.3}\label{eq:15.1.3}\\
  \Gamma^i{}_{tj}
  &=\frac{\dot a}{a}\delta^i_j
  =H\delta^i_j,
  \tag{15.1.4}\label{eq:15.1.4}\\
  \Gamma^i{}_{jk}
  &=\frac12\widetilde g^{i\ell}
    \left(
      \partial_j\widetilde g_{\ell k}
      +\partial_k\widetilde g_{\ell j}
      -\partial_\ell\widetilde g_{jk}
    \right)
  \equiv\widetilde\Gamma^i{}_{jk}.
  \tag{15.1.5}\label{eq:15.1.5}
\end{align}
Its Ricci tensor then has the elements
% MODERNIZATION CHECK: (15.1.6), (15.1.8), and the spatial Ricci
% definition (15.1.9) have been reversed from the source curvature sign.
\begin{align}
  R_{tt}
  &=-\frac{3\ddot a}{a},
  \tag{15.1.6}\label{eq:15.1.6}\\
  R_{ti}
  &=0,
  \tag{15.1.7}\label{eq:15.1.7}\\
  R_{ij}
  &=\widetilde R_{ij}
    +(a\ddot a+2\dot a^2)\widetilde g_{ij},
  \tag{15.1.8}\label{eq:15.1.8}
\end{align}
where \(\widetilde R_{ij}\) is the spatial Ricci tensor calculated from the
metric \(\widetilde g_{ij}\):
\begin{equation}
  \widetilde R_{ij}
  =
  \partial_k\widetilde\Gamma^k{}_{ij}
  -\partial_j\widetilde\Gamma^k{}_{ki}
  +\widetilde\Gamma^k{}_{k\ell}\widetilde\Gamma^\ell{}_{ij}
  -\widetilde\Gamma^k{}_{j\ell}\widetilde\Gamma^\ell{}_{ki}.
  \tag{15.1.9}\label{eq:15.1.9}
\end{equation}
Instead of calculating \(\widetilde R_{ij}\) directly, we can save ourselves
a good deal of work by recalling that \(\widetilde g_{ij}\), as the metric
of a maximally symmetric space, must necessarily have a Ricci tensor of the
form \eqref{eq:13.2.4}:
\begin{equation}
  \widetilde R_{ij}=2k\widetilde g_{ij}.
  \tag{15.1.10}\label{eq:15.1.10}
\end{equation}
Together with \eqref{eq:15.1.8}, this gives the space-space components of the
spacetime Ricci tensor as
\begin{equation}
  R_{ij}
  =(a\ddot a+2\dot a^2+2k)\widetilde g_{ij}.
  \tag{15.1.11}\label{eq:15.1.11}
\end{equation}

As shown in Section~\ref{sec:14.2}, the energy-momentum tensor here must have
the perfect-fluid form
\begin{equation}
  T_{\mu\nu}
  =p\,g_{\mu\nu}+(\rho+p)U_\mu U_\nu,
  \tag{15.1.12}\label{eq:15.1.12}
\end{equation}
where \(p\) and \(\rho\) are functions of \(t\) alone, and \(U^\mu\) is
given by Eqs.~\eqref{eq:14.2.13} and \eqref{eq:14.2.14}:
\begin{equation}
  U^t=1,\qquad U^i=0.
  \tag{15.1.13}\label{eq:15.1.13}
\end{equation}
The source term in the Einstein equations is then
\begin{equation}
  S_{\mu\nu}
  \equiv T_{\mu\nu}-\frac12g_{\mu\nu}T^\lambda{}_\lambda
  =\frac12(\rho-p)g_{\mu\nu}+(\rho+p)U_\mu U_\nu,
  \tag{15.1.14}\label{eq:15.1.14}
\end{equation}
so Eqs.~\eqref{eq:15.1.1}, \eqref{eq:15.1.13}, and
\eqref{eq:15.1.14} give
\begin{align}
  S_{tt}
  &=\frac12(\rho+3p),
  \tag{15.1.15}\label{eq:15.1.15}\\
  S_{it}
  &=0,
  \tag{15.1.16}\label{eq:15.1.16}\\
  S_{ij}
  &=\frac12(\rho-p)a^2\widetilde g_{ij}.
  \tag{15.1.17}\label{eq:15.1.17}
\end{align}
The Einstein equations read
% MODERNIZATION CHECK: the source prints R_{\mu\nu}=-8\pi G S_{\mu\nu};
% the sign is reversed for the binding curvature convention.
\[
  R_{\mu\nu}=8\pi G S_{\mu\nu}.
\]
With Eqs.~\eqref{eq:15.1.6}, \eqref{eq:15.1.7},
\eqref{eq:15.1.11}, and \eqref{eq:15.1.15}--\eqref{eq:15.1.17}, the
time-time component gives
\begin{equation}
  3\ddot a=-4\pi G(\rho+3p)a,
  \tag{15.1.18}\label{eq:15.1.18}
\end{equation}
the space-space components give the single equation
\begin{equation}
  a\ddot a+2\dot a^2+2k
  =4\pi G(\rho-p)a^2,
  \tag{15.1.19}\label{eq:15.1.19}
\end{equation}
and the space-time components give \(0=0\).

By eliminating \(\ddot a\) from Eqs.~\eqref{eq:15.1.18} and
\eqref{eq:15.1.19}, we find a first-order differential equation for \(a(t)\):
\begin{equation}
  \dot a^2+k=\frac{8\pi G}{3}\rho a^2.
  \tag{15.1.20}\label{eq:15.1.20}
\end{equation}
In addition we have the equation of energy conservation
\eqref{eq:14.2.19}, which in modern form is
\[
  \dot\rho+3H(\rho+p)=0,
\]
or, equivalently,
\begin{equation}
  \frac{\dd}{\dd a}\bigl(\rho a^3\bigr)=-3pa^2.
  \tag{15.1.21}\label{eq:15.1.21}
\end{equation}
Given an equation of state \(p=p(\rho)\), we can use this equation to
determine \(\rho\) as a function of \(a\). For instance, if the energy
density of the universe is dominated by nonrelativistic matter with negligible
pressure, then Eq.~\eqref{eq:15.1.21} gives
\begin{equation}
  \rho\propto a^{-3}
  \qquad\text{for }p\ll\rho,
  \tag{15.1.22}\label{eq:15.1.22}
\end{equation}
whereas if the energy density is dominated by relativistic particles, such as
photons, then \(p=\rho/3\), and Eq.~\eqref{eq:15.1.21} gives
\begin{equation}
  \rho\propto a^{-4}
  \qquad\text{for }p=\frac{\rho}{3}.
  \tag{15.1.23}\label{eq:15.1.23}
\end{equation}

Knowing \(\rho\) as a function of \(a\), we can determine \(a(t)\) for all
time by solving Eq.~\eqref{eq:15.1.20}. \emph{The fundamental equations of
dynamical cosmology are thus the Einstein equations\allowbreak
\eqref{eq:15.1.18}--\eqref{eq:15.1.20}, the energy-conservation equation
\eqref{eq:15.1.21}, and the equation of state.} The cosmological models,
based on a Robertson--Walker metric, in which \(a(t)\) is derived in this way,
are known as Friedmann models.\hyperref[ch15-ref-1]{\textsuperscript{1}}

[Incidentally, the solution \(a(t)\) determined in this way will
automatically satisfy Eqs.~\allowbreak\eqref{eq:15.1.18}
and\allowbreak\ \eqref{eq:15.1.19}, for by differentiating
Eq.~\allowbreak\eqref{eq:15.1.20} with respect to time and using
Eq.~\eqref{eq:15.1.21}, we find
\begin{align*}
  2\dot a\ddot a
  &=
  \frac{8\pi G\dot a}{3a}
  \left[
    -\rho a^2+\frac{\dd}{\dd a}(\rho a^3)
  \right]\\
  &=
  \frac{8\pi G\dot a}{3a}
  \left(-\rho a^2-3pa^2\right),
\end{align*}
which is just the same as Eq.~\eqref{eq:15.1.18}.
Equation~\eqref{eq:15.1.19} then follows trivially from
Eqs.~\eqref{eq:15.1.18} and \eqref{eq:15.1.20}. The reason we can make
do with the single field equation \eqref{eq:15.1.20}, instead of the two
equations \eqref{eq:15.1.18} and \eqref{eq:15.1.19}, is of course that
these two equations are not functionally independent, being related by the
Bianchi identities to the energy-conservation equation
\eqref{eq:15.1.21}.]

It is possible to learn a good deal about the past and future expansion of the
universe by simply inspecting Eqs.~\eqref{eq:15.1.18}--\eqref{eq:15.1.21},
even without specifying a definite equation of state.
Equation~\eqref{eq:15.1.18} shows that as long as the quantity
\(\rho+3p\) remains positive, the ``acceleration'' \(\ddot a/a\) is
negative. Since at present \(a>0\) (by definition), and \(\dot a/a>0\)
(because we see red shifts, not blue shifts), it follows that the curve of
\(a(t)\) versus \(t\) must be concave downward, and \emph{must have
reached \(a(t)=0\) at some finite time in the past}. Let us call this time
\(t=0\), so that
\begin{equation}
  a(0)=0.
  \tag{15.1.24}\label{eq:15.1.24}
\end{equation}
The present time \(t_0\) is then the time elapsed since this singularity, and
may justly be called the age of the universe. If \(\ddot a(t)\) vanished for
\(0<t<t_0\), then \(a(t)\) would be just \(a(t_0)t/t_0\), and so the
age \(t_0\) would just equal the Hubble time
\(H_0^{-1}=a(t_0)/\dot a(t_0)\). With \(\ddot a\) negative for
\(0<t<t_0\), \emph{the age of the universe must be less than the Hubble
time}:
\begin{equation}
  t_0<H_0^{-1}.
  \tag{15.1.25}\label{eq:15.1.25}
\end{equation}

Looking into the future, we see from Eq.~\eqref{eq:15.1.21} that as long as
the pressure \(p\) does not become negative, the density \(\rho\) must
decrease with increasing \(a\) at least as fast as \(a^{-3}\), so that for
\(a\to\infty\), the right-hand side of Eq.~\eqref{eq:15.1.20} vanishes at
least as fast as \(a^{-1}\). For \(k=-1\), \(\dot a^2(t)\) remains
positive-definite, so \(a(t)\) goes on increasing, with
\[
  a(t)\longrightarrow t
  \qquad\text{as }t\longrightarrow\infty
  \qquad\text{for }k=-1.
\]
For \(k=0\), \(\dot a^2(t)\) remains positive-definite, so \(a(t)\) goes
on increasing, but more slowly than \(t\). For \(k=+1\), \(\dot a^2(t)\)
will reach zero when \(\rho a^2\) drops to the value \(3/(8\pi G)\).
Since \(\ddot a\) is negative-definite, \(a(t)\) will then begin to
decrease again, and eventually must again reach \(a=0\) at some finite time
in the future. Hence the qualitative course of cosmic history is determined
by the sign of the spatial curvature: \emph{If \(k=-1\) or \(k=0\), then
the universe will go on expanding forever, whereas if \(k=+1\), then the
expansion will eventually cease and be followed by a contraction back to a
singular state with \(a(t)=0\).}

The combination of the Cosmological Principle with the Einstein field
equations illuminates some of the profound questions raised by Newton and Mach.
(See Section~\ref{sec:1.3}.) Suppose that we want to study some physical
system \(\mathcal S\), such as the solar system or the rotating bucket of
Newton, whose size is much less than the cosmic scale factor \(a\). We can
imagine \(\mathcal S\) to be placed in a spherical cavity, cut out of the
expanding universe, and so long as the size of this cavity is much less than
\(a\), we can safely consider this cavity to be empty apart from the system
\(\mathcal S\). If \(\mathcal S\) were absent, the gravitational field
inside the cavity would be a spherically symmetric field with
\(R_{\mu\nu}=0\), and hence, according to the Birkhoff theorem (see
Section~\ref{sec:11.7}), it would have a flat-space metric equivalent to the
Minkowski metric \(\eta_{\mu\nu}\). As long as the system \(\mathcal S\) is
not too big, we can then calculate its gravitational field as a perturbation on
\(\eta_{\mu\nu}\), ignoring all matter outside our cavity, and we can
determine the behavior of the system by using Newtonian or special-relativistic
mechanics. The question of what determines the inertial frames is now answered,
for the only reference frames in which the whole universe appears spherically
symmetric, so that Birkhoff's theorem applies, are the frames at the center of
our cavity, which do not rotate with respect to the expanding cloud of
``typical'' galaxies. The inertial frames are any reference frames that move at
constant velocity, and without rotation, relative to the frames in which the
universe appears spherically symmetric.

These remarks lead to an alternative
derivation\hyperref[ch15-ref-2]{\textsuperscript{2}} of the dynamical equations
for an expanding universe. If we mentally draw a comoving spherical surface
anywhere in the universe, then as long as its proper radius is much less than
\(a(t)\), the galaxies within this sphere will move under the influence of
their own gravitational field, and the gravitational field of the rest of the
universe may be neglected. We can then think of the universe as consisting of
a Newtonian gas in a state of everywhere-uniform expansion. Any given gas
particle will have a trajectory
\[
  \mathbf{x}(t)
  =\mathbf{x}(t_0)\frac{a(t)}{a(t_0)},
\]
with \(a(t)\) a scale factor common to the whole gas. [Note that the gas
appears the same to an observer mounted on any gas particle as it does to an
observer at the origin. Also, the ``comoving'' coordinates of a given gas
particle are not \(x^i(t)\), but rather
\(r^i=x^i(t)/a(t)\).] The gravitational potential energy \(V\) of such a
particle just arises from the matter within a sphere of radius
\(\lvert\mathbf{x}(t)\rvert\) and center at the origin, so
\begin{align*}
  V(t)
  &=-\frac{4\pi}{3}\lvert\mathbf{x}(t)\rvert^3\rho(t)
    \frac{mG}{\lvert\mathbf{x}(t)\rvert}\\
  &=-\frac{4\pi}{3}mG\lvert\mathbf{x}(t_0)\rvert^2\rho(t)
    \frac{a^2(t)}{a^2(t_0)},
\end{align*}
where \(m\) is the particle mass and \(\rho(t)\) is the uniform mass
density of the gas. The kinetic energy of this particle is
\[
  T(t)
  =\frac12m\lvert\dot{\mathbf{x}}(t)\rvert^2
  =\frac12m\lvert\mathbf{x}(t_0)\rvert^2
    \frac{\dot a^2(t)}{a^2(t_0)},
\]
and its total energy is
\[
  E\equiv T(t)+V(t)
  =\frac12m\frac{\lvert\mathbf{x}(t_0)\rvert^2}{a^2(t_0)}
    \left[
      \dot a^2(t)-\frac{8\pi G}{3}\rho(t)a^2(t)
    \right].
\]
With \(E\) constant, this is just the same as Eq.~\eqref{eq:15.1.20},
provided that we identify the energy of a particle as
\begin{equation}
  E
  =-\frac12m\frac{\lvert\mathbf{x}(t_0)\rvert^2}{a^2(t_0)}\,k.
  \tag{15.1.26}\label{eq:15.1.26}
\end{equation}
For \(k=-1\), \(E\) is positive-definite, so gravitation cannot prevent
the gas from dispersing to infinity, with a finite asymptotic velocity. For
\(k=0\), \(E\) vanishes, and the gas is just barely able to expand
indefinitely. For \(k=+1\), \(E\) is negative, and the explosion must
ultimately cease and be followed by an implosion.

Although Newtonian cosmology can reproduce the chief results derived from
Einstein's equations, it is essentially incomplete, for several reasons. We
need general relativity to justify the neglect of all the matter outside a sphere
of radius \(\lvert\mathbf{x}(t)\rvert\) in calculating the gravitational
potential at \(\mathbf{x}(t)\). We cannot use Newtonian mechanics when the
medium itself consists of particles with relativistic local velocities.
Finally, it is only through the use of general relativity that we are able to
interpret the observation of light signals correctly in terms of the cosmic
scale factor \(a(t)\).
